\documentclass[aspectratio=169,11pt]{beamer}
\usetheme{default}
\setbeamertemplate{navigation symbols}{}
\usepackage{graphicx}
\usepackage[T1]{fontenc}
\definecolor{ncsured}{HTML}{CC0000}
\definecolor{ncsudk}{HTML}{990000}
\setbeamercolor{structure}{fg=ncsured}
\setbeamercolor{frametitle}{fg=white,bg=ncsured}
\setbeamercolor{title}{fg=ncsudk}
\setbeamercolor{itemize item}{fg=ncsured}
\setbeamercolor{itemize subitem}{fg=ncsured}
\setbeamerfont{frametitle}{series=\bfseries}
\setbeamertemplate{frametitle}{%
  \nointerlineskip\vskip2pt%
  \begin{beamercolorbox}[wd=\paperwidth,ht=2.4ex,dp=1ex,leftskip=0.3cm]{frametitle}%
  \bfseries\insertframetitle%
  \end{beamercolorbox}}
\setbeamertemplate{footline}{%
  \hbox{\begin{beamercolorbox}[wd=\paperwidth,ht=2.2ex,dp=1ex,leftskip=.3cm,rightskip=.3cm]{}%
  \tiny Optimal Market Making for Cryptocurrency Options \hfill Week 4 \hfill \insertframenumber%
  \end{beamercolorbox}}}

\title{Week 4: Stochastic Optimal Control in Continuous Time}
\subtitle{The general theory behind the worked example}
\date{North Carolina State University \,\textbar\, Summer 2026}
\titlegraphic{\includegraphics[width=4.2cm]{/Users/dendisuhubdy/Github/ncsu/assets/ncsu_logo.png}}
\begin{document}
\begin{frame}[plain]\titlepage
\begin{center}\tiny Industry Mentor: Dendi Suhubdy (Bitwyre)  ·  Guest: Fenni Kang (Coincall)\end{center}\end{frame}
\begin{frame}{Learning Objectives}
\begin{itemize}
  \item Apply the dynamic programming principle in continuous time
  \item Understand the HJB equation and the verification theorem
  \item See why CARA and CRRA utilities give tractable problems
\end{itemize}\end{frame}
\begin{frame}{Lecture 1: Dynamic Programming \& HJB}
\begin{itemize}
  \item Controlled diffusions and admissible controls
  \item The dynamic programming principle
  \item The infinitesimal generator
  \item Deriving the HJB equation; boundary/terminal conditions
\end{itemize}\end{frame}
\begin{frame}{Lecture 2: Verification \& Utilities}
\begin{itemize}
  \item The verification theorem: a smooth HJB solution IS the value function
  \item Viscosity solutions: when value functions are non-smooth
  \item CARA vs. CRRA vs. HARA
  \item Why CARA decouples the value function from wealth
\end{itemize}\end{frame}
\begin{frame}{Recitation}
\begin{itemize}
  \item Solve the Merton problem under CARA, then CRRA
  \item Compare the two solutions side by side
  \item Recognize the recurring HJB -\textgreater{} ansatz -\textgreater{} ODE -\textgreater{} solution pattern
\end{itemize}\end{frame}
\begin{frame}{Reading \& Problem Set 4 (Python/PyTorch)}
\begin{itemize}
  \item Pham (2009) Ch. 3-4; Cartea-Jaikumar-Penalva Ch. 6
  \item PS4: Merton CRRA/CARA, Almgren-Chriss, HJB-residual check
\end{itemize}\end{frame}
\end{document}